\section{相关工作}

与其他大型分布式文件系统(例如AFS[5])一样，GFS提供位置无关的命名空间，
使数据能够为了负载均衡或容错而被透明地迁移。与AFS不同，为了提供聚合性能并提升容错能力，
GFS以更接近xFS[1]和Swift[3]的方式，将一个文件的数据分布到多个存储服务器上。

由于磁盘相对廉价，并且复制比更复杂的RAID[9]方案更简单，GFS目前只使用复制来提供冗余。
因此，它比xFS或Swift消耗更多原始存储空间。

与AFS、xFS、Frangipani[12]和Intermezzo[6]等系统不同，GFS不在文件系统接
口之下提供任何缓存。我们的目标工作负载在单次应用运行中几乎没有数据复用，
因为它们要么流式处理大型数据集，要么在其中随机寻道，并且每次只读取少量数据。

一些分布式文件系统，例如Frangipani、xFS、明尼苏达大学的GFS[11]和GPFS[10]，
移除了中心化服务器，并依赖分布式算法实现一致性与管理。我们选择中心化方案，以简化设计、
提高可靠性并获得灵活性。特别是，中心化master使实现复杂的chunk放置与复制策略容易得多，
因为master已经拥有大部分相关信息，并控制这些信息如何变化。我们通过保持master状态较小，
并将其完整复制到其他机器上来实现容错。目前，可扩展性与读取高可用性由我们的
shadow master机制提供。对master状态的更新通过追加到预写日志而持久化。因此，
我们可以采用类似Harp[7]的primary-copy方案，提供比当前方案具有更强一致性保证的高可用性。

在向大量客户端提供聚合性能方面，我们处理的问题与Lustre[8]类似。
不过，我们专注于应用需求，而非构建符合POSIX的文件系统，因此显著简化了这个问题。
此外，GFS假设系统由大量不可靠组件构成，因此容错是设计的核心。

GFS与NASD架构[4]最为相似。NASD架构以网络附加磁盘驱动器为基础，而GFS如NASD
原型一样，使用普通机器作为chunkserver。与NASD工作不同，我们的chunkserver
使用延迟分配的固定大小chunk，而非可变长度对象。此外，GFS实现了负载再平衡、
复制与恢复等生产环境所需的功能。

与明尼苏达大学的GFS和NASD不同，我们并不试图改变存储设备模型。
我们关注利用现有普通组件，满足复杂分布式系统日常的数据处理需求。

原子\textbf{record append}支持的生产者-消费者队列，
解决了与River[2]中分布式队列类似的问题。River使用分布在多台机器上的内存队列
和精细的数据流控制，而GFS使用可由多个生产者并发追加的持久化文件。
River模型支持m对n的分布式队列，但缺少持久化存储带来的容错能力；
GFS则只能高效支持m对1的队列。多个消费者可以读取同一文件，但必须协调对输入负载的划分。
